\documentclass[11pt,a4paper]{article}

% ============================================================================
% PACKAGES (matching the GAUSE method deep-dive)
% ============================================================================
\usepackage[utf8]{inputenc}
\usepackage[T1]{fontenc}
\usepackage{lmodern}
\usepackage{amsmath,amssymb,amsthm}
\usepackage{mathtools}
\usepackage{bm}
\usepackage{algorithm}
\usepackage{algorithmic}
\usepackage{booktabs}
\usepackage{multirow}
\usepackage{array}
\usepackage{float}
\usepackage{graphicx}
\usepackage{xcolor}
\usepackage[numbers]{natbib}
\usepackage{hyperref}
\usepackage{tikz}
\usetikzlibrary{shapes,arrows,positioning,calc,patterns,decorations.pathreplacing,shapes.geometric,arrows.meta,fit,backgrounds,intersections}
\usepackage{colortbl}
\usepackage{pgfplots}
\pgfplotsset{compat=1.17}
\usepgfplotslibrary{fillbetween}
\usepackage{listings}
\usepackage{tcolorbox}
\usepackage{enumitem}
\usepackage{geometry}
\geometry{margin=1in}

% ============================================================================
% THEOREM ENVIRONMENTS
% ============================================================================
\theoremstyle{definition}
\newtheorem{definition}{Definition}[section]
\newtheorem{example}{Example}[section]
\newtheorem{remark}{Remark}[section]

\theoremstyle{plain}
\newtheorem{theorem}{Theorem}[section]
\newtheorem{proposition}{Proposition}[section]
\newtheorem{lemma}{Lemma}[section]
\newtheorem{corollary}{Corollary}[section]

% ============================================================================
% CUSTOM COMMANDS
% ============================================================================
\newcommand{\R}{\mathbb{R}}
\newcommand{\E}{\mathbb{E}}
\newcommand{\Var}{\mathrm{Var}}
\newcommand{\Cov}{\mathrm{Cov}}
\newcommand{\N}{\mathcal{N}}
\newcommand{\KL}{\mathrm{KL}}
\newcommand{\TV}{\mathrm{TV}}
\newcommand{\nm}{\textsc{RISP}}
\newcommand{\Ginv}{G_{\mathrm{inv}}}
\newcommand{\Gforget}{G_{\mathrm{forget}}}
\newcommand{\argmax}{\operatornamewithlimits{arg\,max}}
\newcommand{\argmin}{\operatornamewithlimits{arg\,min}}

% Code listing style
\lstset{
    language=Python,
    basicstyle=\ttfamily\small,
    keywordstyle=\color{blue},
    commentstyle=\color{gray},
    stringstyle=\color{red},
    numbers=left,
    numberstyle=\tiny\color{gray},
    frame=single,
    breaklines=true,
    captionpos=b
}

% Colored boxes for intuition
\newtcolorbox{intuition}[1][]{
    colback=blue!5!white,
    colframe=blue!75!black,
    fonttitle=\bfseries,
    title=Intuition,
    #1
}

\newtcolorbox{keypoint}[1][]{
    colback=green!5!white,
    colframe=green!50!black,
    fonttitle=\bfseries,
    title=Key Point,
    #1
}

\newtcolorbox{warning}[1][]{
    colback=red!5!white,
    colframe=red!75!black,
    fonttitle=\bfseries,
    title=Important,
    #1
}

% ============================================================================
% DOCUMENT
% ============================================================================
\title{\Huge\textbf{\nm{}: A Mathematical Deep Dive}\\[0.5cm]
\Large Retention and Invariance in Decision-Focused\\
Strategy Pools for Non-Stationary Markets}

\author{Yuhao Li\\
University of Pennsylvania\\
\texttt{li88@sas.upenn.edu}}

\date{June 2026}

\begin{document}

\maketitle

\begin{abstract}
This document is the exhaustive, ground-up mathematical companion to the
\nm{} research paper (\emph{When the Regime Returns}) and its explanatory
report. It develops, from first principles, every tool the framework's
theory uses---decision-focused learning and the SPO+ surrogate,
concentration of measure, distances between distributions and Pinsker's
inequality, exponentiated-gradient dynamics, and invariant risk across
environments---then assembles them into the four results that organize the
contribution: the episode-transfer lemma (the invariance gap), the
eviction-rate lemma (the forgetting deficit), the post-reactivation regret
decomposition with its predicted-and-confirmed super-additive interaction,
and the break-even-dormancy pricing of retention against carrying costs.
The framework's central claim is that financial non-stationarity is
two-layered---regimes \emph{recur} (so dormant expertise is worth
retaining, and any reward-driven capacity allocator provably fails to
retain it) and recurrences \emph{differ} (so retained expertise must be
trained for across-episode invariance, or it faithfully preserves the last
episode's accidents)---and that the two layers are governed by independent
design choices whose product, but not either factor alone, recovers a
hand-designed oracle assignment without hand design. Every theoretical
statement is given with its full proof or an explicitly labeled sketch;
every assumption carries an audit of its idealizations; every experimental
number in the evidence part traces to a released result file; and the
document reports, at full prominence, the three places where our
pre-registered expectations were refuted by our own experiments (the
real-data structure gate, the competition-batching story, and the
high-signal regime where the entire mechanism inverts). \textbf{Prerequisites:}
probability at the level of a first graduate course, linear algebra, and
convexity basics; everything else is developed here.
\end{abstract}

\tableofcontents
\newpage

